\section{Essential Tools}

Before you read any farther, you need to know about the \textbf{tools} that are \textbf{needed} to succeed in this course. You will need access to the \textbf{textbook}, a \textbf{computer} with the software listed below, and the \textbf{time} to learn on your own and complete assignments.

\subsection{Your Textbook}

This web-book document that is available on Moodle \textbf{is not} the textbook, it is a guide to the course with readings, assignments, documents, etc. collected together. The textbook for this course is\ldots

\begin{quote}
\textit{Modern Physical Organic Chemistry} by Eric V. Anslyn and Dennis A. Dougherty. MIT University Press.
\end{quote}

\textbf{This book is a tour de force.} I encourage you to buy the textbook and \textbf{keep it} for life. It covers the basics of physical organic chemistry and then has chapters that \textbf{can be applied in just about any field of chemistry}. Supramolecular chemistry? Check. Conductive polymers? Check. Theoretical chemistry? Check.

The textbook is available in the bookstore and other options for textbook access are described above in the course information section. \textbf{All chapters and problem assignments will refer to this book.}

The textbook for this course was \textbf{written} by two eminent physical organic chemists and then \textbf{peer reviewed} and corrected by dozens more. Look around you. How many students are in this class? It's not 200. I guarantee that \textbf{no one is making money} from this book. If we do not support efforts like this, we may never see a second edition.

These authors are far too old for a \emph{Patreon} account. Please \textbf{do not pirate} textbooks. I know that downloading is easy to do, but someday you may be writing a textbook of your own. Start building your \textbf{karma} now.

\subsection{The Literature}

We will be using literature examples and you will be expected to \textbf{become familiar with searching and using chemical literature}. I will try to use \textbf{only papers that are accessible through the UPEI Library} and its electronic subscriptions. If I use a paper that is only available through interlibrary loan, I will make it available through Moodle.

\textbf{Papers that you reference} in reports or assignments \textbf{must be accessible through the UPEI Library system} or by interlibrary loan. If obtained through interlibrary loan, a copy of the paper must be submitted with your work.

\textbf{We will not pirate papers}, regardless of our opinions about scientific publishing and profit-driven publication models.

The \textbf{library} provides a useful search tool that can be used to locate references. If the library does not hold a \textbf{license} for a source, you will usually be able to obtain it through \textbf{interlibrary loan}, provided that you \textbf{give yourself enough time}. Begin research projects early.

Here are some useful literature resources:

\begin{itemize}
\item \href{https://library.upei.ca/scifinder}{ACS SciFinder via UPEI}
\item \href{https://scholar.google.ca}{Google Scholar}
\item \href{https://library.upei.ca/subject/chemistry}{UPEI Library Chemistry Resource Page}
\end{itemize}

\subsection{Your Computer}

Welcome to the modern world. For the past 20 years, students have \textbf{needed computers to succeed} in university.

You will need \textbf{access} to a computer with the \textbf{tools} listed below. These tools will be \textbf{available} on many \textbf{campus} computers. I have selected only \textbf{free} software options that you can install on your own computer.

\textbf{We will not be programming}. We will be learning how to use software tools. Any coding that appears in the course will be provided in advance, and I likely borrowed it from somewhere on the internet. I am not a coder, I am a \textbf{user}, and \href{https://en.wikipedia.org/wiki/Tron}{users are the true heroes}.

Below is a list of software that will be available on many campus computers or can be installed on your own computer.

\begin{enumerate}
\item \textbf{A chemical drawing program}
    \begin{itemize}
    \item There are many options. I recommend the \emph{Marvin} package from ChemAxon:
    \href{https://chemaxon.com/products/marvin}{ChemAxon Marvin}.
    
    \item It runs on multiple operating systems.
    
    \item A convenient web-based version is available through \emph{Chemicalize}:
    \href{https://chemicalize.com/welcome}{Chemicalize}.
    
    \item The website provides access to many of the features of Marvin without installing software.
    \end{itemize}

\item \textbf{The usual productivity software}
    \begin{itemize}
    \item You will need a \textbf{word processor}, \textbf{spreadsheet}, and \textbf{presentation program}.
    
    \item A free option is
    \href{https://www.libreoffice.org/}{LibreOffice}.
    
    \item \href{https://www.upei.ca/learning2020/technology-basics/useful-applications}{Office 365}
    is also available to all UPEI students.
    
    \item In this case, I recommend the commercial option because it is already available through the university.
    \end{itemize}

\item \textbf{Python with the Jupyter Notebook environment}
    \begin{itemize}
    \item We will use \textbf{Python} for data analysis and plotting.
    
    \item It is a powerful scientific tool that you will likely continue using throughout your career.
    
    \item We will access Python through the \textbf{Google Colab} website, so there is no need to install software locally.
    \end{itemize}
\end{enumerate}

\newpage
\subsection{Your Time}

This course will require your time. Only you can explore the concepts presented in the lessons by \textbf{investing your time}.

You will need:

\begin{itemize}
\item \textbf{time to read} the textbook and consider the suggested problems;
\item \textbf{time to struggle} with software tools as we explore data analysis and computational chemistry;
\item \textbf{time to write} reports and prepare presentations; and
\item \textbf{time to relax}.
\end{itemize}